% GENERATED FILE. Edit canonical lessons, then run npm run generate:books.
% canonical-source-hash: fnv1a64:30bdce3af20d80f4
% canonical-lessons: LA-C207-saeculum, LA-C207-aetas, LA-C207-occasus, LA-C207-ortus, LA-C207-diluculum

\chapter{Generation, Age, Sunset, Sunrise, First light}
\label{ch:la-generation-age-sunset-sunrise-first-light}
\hlchaptermodality{207}

\begin{chapteropening}
\textbf{By the end of this chapter:} \emph{I can say a generation, age, sunset, sunrise and first light.}

Complete the last lesson of chapter 207: I can say a generation, age, sunset, sunrise and first light.
\end{chapteropening}

\section[saeculum, saeculī]{saeculum, saeculī — a generation, a century}

\label{lesson:LA-C207-saeculum}



\begin{quote}
Before the new one: say the Latin for talent, then the Latin for a habit.
\end{quote}

\subsection*{You'll want to know: saeculum, saeculī}
\textbf{saeculum, saeculī} — \textquotedblleft{}a generation, a century\textquotedblright{}.

\begin{cousinweb}
\textbf{saeculum} has no sure origin; Lewis \& Short call it doubtful. Welsh \emph{hoedl}, \textquotedblleft{}lifetime\textquotedblright{}, is sometimes compared. English \emph{secular} comes from it.
\end{cousinweb}

\begin{grammarlens}[title={What you've built}]
One more word for this chapter.
\end{grammarlens}

\subsection*{Your turn}
\begin{itemize}
\raggedright
  \item \emph{Say it:} \emph{saeculum, saeculī}
  \item \emph{Say it:} \emph{saeculum, saeculī}, once more
  \item \emph{Say it:} say \emph{saeculum, saeculī}
  \item \emph{Recall:} say the Latin for talent, then the Latin for a habit, then say \emph{saeculum, saeculī} again
\end{itemize}

\begin{tcolorbox}[breakable,colback=teal!4,colframe=teal!35!black,title={Before you move on}]
What does saeculum, saeculī mean? (\textquotedblleft{}A generation, a century\textquotedblright{}.) Say it once more.
\end{tcolorbox}
\section[aetās, aetātis]{aetās, aetātis — age, a lifetime}

\label{lesson:LA-C207-aetas}



\begin{quote}
Before the new one: say the Latin for a habit, then the Latin for a generation.
\end{quote}

\subsection*{You'll want to know: aetās, aetātis}
\textbf{aetās, aetātis} — \textquotedblleft{}age, a lifetime\textquotedblright{}.

\begin{cousinweb}
\textbf{aetās} is contracted from older \emph{aevitās} (Lewis \& Short), built on \emph{aevum}, \textquotedblleft{}a lifetime, an age\textquotedblright{}, which is usually traced to Proto-Indo-European \emph{*h2eyu-}, \textquotedblleft{}life force\textquotedblright{}.
\end{cousinweb}

\begin{grammarlens}[title={What you've built}]
One more word for this chapter.
\end{grammarlens}

\subsection*{Your turn}
\begin{itemize}
\raggedright
  \item \emph{Say it:} \emph{aetās, aetātis}
  \item \emph{Say it:} \emph{aetās, aetātis}, once more
  \item \emph{Say it:} say \emph{aetās, aetātis}
  \item \emph{Recall:} say the Latin for a habit, then the Latin for a generation, then say \emph{aetās, aetātis} again
\end{itemize}

\begin{tcolorbox}[breakable,colback=teal!4,colframe=teal!35!black,title={Before you move on}]
What does aetās, aetātis mean? (\textquotedblleft{}Age, a lifetime\textquotedblright{}.) Say it once more.
\end{tcolorbox}
\section[occāsus, occāsūs]{occāsus, occāsūs — sunset; a setting}

\label{lesson:LA-C207-occasus}



\begin{quote}
Before the new one: say the Latin for a generation, then the Latin for age.
\end{quote}

\subsection*{You'll want to know: occāsus, occāsūs}
\textbf{occāsus, occāsūs} — \textquotedblleft{}sunset; a setting\textquotedblright{}.

\begin{cousinweb}
\textbf{occāsus} is built on \emph{occidō}, \textquotedblleft{}to fall down, to set\textquotedblright{} (Lewis \& Short), from \emph{ob}, \textquotedblleft{}towards\textquotedblright{}, and \emph{cadō}, \textquotedblleft{}to fall\textquotedblright{}: the sun's going down.
\end{cousinweb}

\begin{grammarlens}[title={What you've built}]
One more word for this chapter.
\end{grammarlens}

\subsection*{Your turn}
\begin{itemize}
\raggedright
  \item \emph{Say it:} \emph{occāsus, occāsūs}
  \item \emph{Say it:} \emph{occāsus, occāsūs}, once more
  \item \emph{Say it:} say \emph{occāsus, occāsūs}
  \item \emph{Recall:} say the Latin for a generation, then the Latin for age, then say \emph{occāsus, occāsūs} again
\end{itemize}

\begin{tcolorbox}[breakable,colback=teal!4,colframe=teal!35!black,title={Before you move on}]
What does occāsus, occāsūs mean? (\textquotedblleft{}Sunset; a setting\textquotedblright{}.) Say it once more.
\end{tcolorbox}
\section[ortus, ortūs]{ortus, ortūs — sunrise; a rising}

\label{lesson:LA-C207-ortus}



\begin{quote}
Before the new one: say the Latin for age, then the Latin for sunset; a setting.
\end{quote}

\subsection*{You'll want to know: ortus, ortūs}
\textbf{ortus, ortūs} — \textquotedblleft{}sunrise; a rising\textquotedblright{}.

\begin{cousinweb}
\textbf{ortus} is built on \emph{orior}, \textquotedblleft{}to rise\textquotedblright{} (Lewis \& Short): the sun's rising. English \emph{orient}, \textquotedblleft{}the east\textquotedblright{}, comes from the same verb.
\end{cousinweb}

\begin{grammarlens}[title={What you've built}]
One more word for this chapter.
\end{grammarlens}

\subsection*{Your turn}
\begin{itemize}
\raggedright
  \item \emph{Say it:} \emph{ortus, ortūs}
  \item \emph{Say it:} \emph{ortus, ortūs}, once more
  \item \emph{Say it:} say \emph{ortus, ortūs}
  \item \emph{Recall:} say the Latin for age, then the Latin for sunset; a setting, then say \emph{ortus, ortūs} again
\end{itemize}

\begin{tcolorbox}[breakable,colback=teal!4,colframe=teal!35!black,title={Before you move on}]
What does ortus, ortūs mean? (\textquotedblleft{}Sunrise; a rising\textquotedblright{}.) Say it once more.
\end{tcolorbox}
\section[dīlūculum, dīlūculī]{dīlūculum, dīlūculī — first light, daybreak}

\label{lesson:LA-C207-diluculum}



\begin{quote}
Before the new one: say the Latin for sunset; a setting, then the Latin for sunrise; a rising.
\end{quote}

\subsection*{You'll want to know: dīlūculum, dīlūculī}
\textbf{dīlūculum, dīlūculī} — \textquotedblleft{}first light, daybreak\textquotedblright{}.

\begin{cousinweb}
\textbf{dīlūculum} is built on \emph{dīlūceō}, \textquotedblleft{}to grow light\textquotedblright{} (Lewis \& Short), from \emph{lūx}, \textquotedblleft{}light\textquotedblright{}, with \emph{-culum}, a diminutive ending: the first glimmer.
\end{cousinweb}

\begin{grammarlens}[title={What you've built}]
One more word for this chapter.
\end{grammarlens}

\subsection*{Your turn}
\begin{itemize}
\raggedright
  \item \emph{Say it:} \emph{dīlūculum, dīlūculī}
  \item \emph{Say it:} \emph{dīlūculum, dīlūculī}, once more
  \item \emph{Say it:} say \emph{dīlūculum, dīlūculī}
  \item \emph{Recall:} say the Latin for sunset; a setting, then the Latin for sunrise; a rising, then say \emph{dīlūculum, dīlūculī} again
\end{itemize}

\begin{tcolorbox}[breakable,colback=teal!4,colframe=teal!35!black,title={Before you move on}]
What does dīlūculum, dīlūculī mean? (\textquotedblleft{}First light, daybreak\textquotedblright{}.) Say it once more.
\end{tcolorbox}
